With both success sets restricted to the common finite-binary recursive semi-Markovian signed-local-monotonicity domain of Definition \(\mathrm{def:model\text{-}class}\), the inclusion
\[
 \operatorname{Succ}(\mathsf{MID})
 \subseteq
 \operatorname{Succ}(\mathsf{MCTFID})
\]
is strict.  A witness is the three-vertex DAG \(G\) with directed edges \(X\to Y\leftarrow Z\), no bidirected edges, positive signs on both directed edges, the observational menu \(\mathcal Z=\{\varnothing\}\), and the positive-denominator query
\[
 Q=P\bigl(Y_{11}=1\mid Y_{10}=1,\,Y_{01}=1\bigr),
\]
where the two subscript coordinates give the interventions on \((X,Z)\).  For this witness \(\operatorname{GID}_{G,M,\mathcal Z}(Q)=1\), so \(\mathsf{MCTFID}\) returns an identifying expression \(F\) satisfying
\[
 F(P_{\mathcal Z}(\mathfrak m))=1
\]
for every compatible model \(\mathfrak m\in\mathfrak M(G,M)\) with \(P_{\mathfrak m}(B_Q)>0\), whereas the published \(\mathsf{MID}\) procedure returns \(\mathsf{FAIL}\).  Hence the published \(\mathsf{MID}\) procedure is not success-set maximal even on the common three-binary-variable unconfounded DAG subclass.  This conclusion does not claim that \(\mathsf{MCTFID}\) solves the broader ordered-variable or general-SCM completeness problem considered by the published \(\mathsf{MID}\) work.